\documentclass[12pt]{article}
\usepackage{amsmath, amssymb}
\usepackage[margin=1in]{geometry}
\usepackage{xcolor}
\usepackage{tcolorbox}

\title{\textbf{Why Does $\sin(2t) + \sin(5t)$ Have Peaks of Amplitude 2?}}
\author{Blackboard Derivation}
\date{}

\begin{document}
\maketitle

\section*{Step 1: The Sum-to-Product Identity}

We begin with the standard trigonometric identity:
\[
\boxed{\sin A + \sin B \;=\; 2\,\cos\!\left(\frac{A-B}{2}\right)\sin\!\left(\frac{A+B}{2}\right)}
\]

\section*{Step 2: Apply to Our Waves}

Let $A = 2t$ and $B = 5t$. Substituting:
\begin{align*}
\sin(2t) + \sin(5t)
&= 2\,\cos\!\left(\frac{2t - 5t}{2}\right)\sin\!\left(\frac{2t + 5t}{2}\right) \\[8pt]
&= 2\,\cos\!\left(\frac{-3t}{2}\right)\sin\!\left(\frac{7t}{2}\right) \\[8pt]
&= 2\,\cos\!\left(\frac{3t}{2}\right)\sin\!\left(\frac{7t}{2}\right)
\end{align*}

where the last step uses $\cos(-x) = \cos(x)$.

\begin{tcolorbox}[colback=blue!5, colframe=blue!50!black, title=Decomposition]
\[
\sin(2t) + \sin(5t) = \underbrace{2\,\cos\!\left(\frac{3t}{2}\right)}_{\text{slow envelope}} \cdot \underbrace{\sin\!\left(\frac{7t}{2}\right)}_{\text{fast carrier}}
\]
\end{tcolorbox}

\section*{Step 3: Finding the Peak Amplitude}

The maximum of the product form occurs when both factors achieve their extreme values simultaneously:
\begin{align*}
\left|\cos\!\left(\frac{3t}{2}\right)\right| &\leq 1 \\[4pt]
\left|\sin\!\left(\frac{7t}{2}\right)\right| &\leq 1
\end{align*}

Therefore:
\[
\left|2\,\cos\!\left(\frac{3t}{2}\right)\sin\!\left(\frac{7t}{2}\right)\right| \;\leq\; 2 \cdot 1 \cdot 1 \;=\; 2
\]

The peak of amplitude $\mathbf{2}$ is achieved when:
\[
\cos\!\left(\frac{3t}{2}\right) = \pm 1 \quad\text{and}\quad \sin\!\left(\frac{7t}{2}\right) = \pm 1 \quad\text{(with matching signs)}
\]

\section*{Step 4: When Exactly Does This Happen?}

\begin{align*}
\cos\!\left(\frac{3t}{2}\right) = 1 &\implies \frac{3t}{2} = 2n\pi \implies t = \frac{4n\pi}{3}, \quad n \in \mathbb{Z} \\[6pt]
\sin\!\left(\frac{7t}{2}\right) = 1 &\implies \frac{7t}{2} = \frac{\pi}{2} + 2m\pi \implies t = \frac{\pi(4m+1)}{7}, \quad m \in \mathbb{Z}
\end{align*}

A peak of exactly $+2$ requires both conditions simultaneously:
\[
\frac{4n\pi}{3} = \frac{\pi(4m+1)}{7} \implies 28n = 3(4m+1) \implies 28n = 12m + 3
\]

This is a Diophantine equation. Since $\gcd(28, 12) = 4$ and $4 \nmid 3$, there is \emph{no} integer solution. The amplitude approaches 2 but never exactly reaches it in finite time --- though it gets arbitrarily close.

\begin{tcolorbox}[colback=green!5, colframe=green!50!black, title=Physical Intuition]
\textbf{Constructive interference:} Both waves crest together $\Rightarrow$ amplitudes add $\Rightarrow$ $1 + 1 = 2$.

\medskip

The bound of $2$ is the \emph{supremum} --- the theoretical maximum that the sum approaches asymptotically as the envelope and carrier nearly align.
\end{tcolorbox}

\end{document}
